% SEGMENT OLP-0040-S01
% Part:sets-functions-relations
% Chapter: size-of-sets
% Section: reduction-alt

\documentclass[../../../include/open-logic-section]{subfiles}

\begin{document}

\olfileid{sfr}{siz}{red-alt}

\olsection{குறைத்தல்}

\begin{editorial}
  இந்தப் பிரிவு \olref[nen-alt]{sec} இன் முடிவுகளுக்கு ஏற்ப,
  குறைத்தல் மூலம் எண்ணிடத்தகாத தன்மையை நிறுவுகிறது.
  \olref[nen]{sec} இன் முடிவுகளுக்கு ஏற்ப அமைந்த,
  சற்றே விரிவான மாற்று வடிவம் \olref[red]{sec} இல் தரப்பட்டுள்ளது.
\end{editorial}

% SEGMENT OLP-0040-S02
$\Bin^\omega$ என்பது !!{nonenumerable} என்பதை
ஒரு மூலைவிட்டமாக்க வாதத்தின் மூலம் நிறுவினோம்.
$\Pow{\Nat}$ என்பது !!{nonenumerable} எனக் காட்டவும்
இதே போன்ற ஒரு மூலைவிட்டமாக்க வாதத்தைப் பயன்படுத்தினோம்.
ஆனால் $\Pow{\Nat}$ என்பது !!{nonenumerable} என நிறுவுவதற்கு
மற்றொரு வழியும் உள்ளது:
\emph{$\Pow{\Nat}$ என்பது !!{enumerable} எனில்,
$\Bin^\omega$ உம் !!{enumerable}} எனக் காட்டலாம்.
$\Bin^\omega$ என்பது !!{nonenumerable} என்று ஏற்கெனவே
அறிந்திருப்பதால், $\Pow{\Nat}$ உம் அவ்வாறே எனப் பின்தொடரும்.

இது ஒரு சிக்கலை மற்றொரு சிக்கலாகக் \emph{குறைத்தல்} எனப்படும்.
இந்த நிலையில், $\Bin^\omega$ ஐப் பட்டியலாக்கும் சிக்கலை,
$\Pow{\Nat}$ ஐப் பட்டியலாக்கும் சிக்கலாகக் குறைக்கிறோம்.
பின்னைய சிக்கலுக்கான ஒரு தீர்வு---$\Pow{\Nat}$ இன் பட்டியலாக்கம்---
முன்னைய சிக்கலுக்கான ஒரு தீர்வை---$\Bin^\omega$ இன்
பட்டியலாக்கத்தை---வழங்கும்.

ஒரு கணம் $B$ ஐப் பட்டியலாக்கும் சிக்கலை,
ஒரு கணம் $A$ ஐப் பட்டியலாக்கும் சிக்கலாகக் குறைப்பதற்கு,
$A$ இன் ஒரு பட்டியலாக்கத்தை $B$ இன் பட்டியலாக்கமாக மாற்றும்
ஒரு வழியைத் தருகிறோம்.
இதனைச் செய்வதற்கான எளிய வழி,
!!a{surjection} $f\colon A \to B$ ஐ வரையறுப்பதாகும்.
$x_1$, $x_2$, \dots{} என்பது $A$ ஐப் பட்டியலாக்கினால்,
$f(x_1)$, $f(x_2)$, \dots{} என்பது $B$ ஐப் பட்டியலாக்கும்.
நமது நிலையில், !!a{surjection} $f\colon \Pow{\Nat}
\to \Bin^\omega$ ஐத் தேடுகிறோம்.

% SEGMENT OLP-0040-S03
\begin{prob}
!!{injective} பண்புடைய சார்பு $g\colon B \to A$ இருந்து,
$B$ என்பது !!{nonenumerable} எனில்,
$A$ உம் அவ்வாறே எனக் காட்டுக.
$A$ இன் ஒரு பட்டியலாக்கத்தை $B$ இன் பட்டியலாக்கமாக மாற்ற
$g$ ஐ எவ்வாறு பயன்படுத்தலாம் எனக் காட்டி இதனைச் செய்க.
\end{prob}

% SEGMENT OLP-0040-S04
\begin{proof}[{\olref[nen-alt]{thm:nonenum-pownat}} ஐக் குறைத்தல் மூலம் நிறுவுதல்]
ஒரு குறைத்தலுக்காக, $\Pow{\Nat}$ என்பது !!{enumerable} எனவும்,
எனவே அதற்கு $N_{1}$, $N_{2}$, $N_{3}$, \dots என்ற
ஒரு பட்டியலாக்கம் உள்ளது எனவும் கொள்க.

% SOURCE CORRECTION OLSIZ-010; see translation/size-notes.tex.
$f(N)$ என்பது பின்வரும் நிபந்தனையை நிறைவேற்றும் $s$ என்ற
தொடராக இருக்குமாறு $f \colon \Pow{\Nat} \to \Bin^\omega$
என்ற சார்பை வரையறுக்கவும்:
$s(n) = 1$ ஆக இருக்கும்போது, அப்போதும் அப்போது மட்டுமே,
$n \in N$; இல்லையெனில் $s(n) = 0$.

இது தெளிவாக ஒரு சார்பை வரையறுக்கிறது;
ஏனெனில் $N \subseteq \Nat$ என இருக்கும்போதெல்லாம்,
எந்த $n \in \Nat$ உம் $N$ இன் !!a{element} ஆக இருக்கும் அல்லது இருக்காது.
எடுத்துக்காட்டாக, இரட்டை இயல் எண்களின் கணம்
$2\Nat = \Setabs{2n}{n \in \Nat} = \{0,2, 4, 6,
\dots\}$ ஆனது $1010101\dots$ என்ற தொடருக்கு அனுப்பப்படுகிறது;
$\emptyset$ என்பது $0000\dots$ க்கும்,
$\Nat$ என்பது $1111\dots$ க்கும் அனுப்பப்படுகின்றன.

இது !!{surjective} பண்பையும் கொண்டுள்ளது:
$0$s, $1$s ஆகியவற்றாலான ஒவ்வொரு தொடரும் ஏதோ ஓர் இயல் எண் கணத்துடன்
ஒத்துள்ளது; அதாவது, தொடரில் $1$s உள்ள இடங்களுக்கு ஒத்த இயல் எண்களை
உறுப்புகளாகக் கொண்ட கணத்துடன் ஒத்துள்ளது.
மேலும் துல்லியமாக, $s \in \Bin^\omega$ எனில்,
$N \subseteq \Nat$ என்பதைப் பின்வருமாறு வரையறுக்கவும்:
\[
N = \Setabs{n \in \Nat}{s(n) = 1}
\]
அப்போது $f$ இன் வரையறையைப் பார்த்து உறுதிசெய்யக்கூடியவாறு,
$f(N) = s$.

இப்போது பின்வரும் பட்டியலைக் கருதுக:
\[
f(N_1), f(N_2), f(N_3), \dots
\]
$f$ என்பது !!{surjective} பண்புடையதால்,
$\Bin^\omega$ இன் ஒவ்வோர் உறுப்பும் ஏதாவது ஓர் உள்ளீட்டிற்கு
$f$ இன் மதிப்பாக இருக்க வேண்டும்; எனவே பட்டியலில் இடம்பெற வேண்டும்.
ஆகவே இந்தப் பட்டியல் $\Bin^\omega$ முழுவதையும் பட்டியலாக்க வேண்டும்.

எனவே $\Pow{\Nat}$ என்பது !!{enumerable} ஆக இருந்தால்,
$\Bin^\omega$ உம் !!{enumerable} ஆக இருக்கும்.
ஆனால் $\Bin^\omega$ என்பது !!{nonenumerable}
(\olref[nen-alt]{thm:nonenum-bin-omega}).
ஆகவே $\Pow{\Nat}$ என்பது !!{nonenumerable}.
\end{proof}

% SEGMENT OLP-0040-S05
%\begin{explain}
%குறைத்தல் எந்தத் திசையில் செல்கிறது என்பதில் குழப்பமடைவது எளிது.
%எடுத்துக்காட்டாக, !!a{surjective} சார்பு $g \colon \Bin^\omega \to X$
%ஒன்று இருப்பது, $X$ என்பது !!{nonenumerable} என நிறுவாது.
%($g(s) = s(1)$ என வரையறுக்கப்படும் $g
%\colon \Bin^\omega \to \Bin$ ஐக் கருதுக;
%இது $0$ களாலும் $1$ களாலும் ஆன ஒரு தொடரை அதன் முதல்
%!!{element} க்கு அனுப்பும் சார்பாகும்.
%சில தொடர்கள் $0$ இலும் சில தொடர்கள் $1$ இலும் தொடங்குவதால்,
%அது மேற்கோர்த்தல் பண்புடையது. ஆனால் $\Bin$ முடிவுறு கணம்.)
%மேலும் $f$ என்பது மேற்கோர்த்தல் பண்புடையதாக இருக்க வேண்டும்;
%இல்லையெனில் வாதம் செயல்படாது:
%$f(x_1)$, $f(x_2)$, \dots{} என்பதில் $Y$ இன் எல்லா
%!!{element}s உம் இடம்பெறும் என்பது அப்போது உறுதியாகாது.
%எடுத்துக்காட்டாக, பின்வருமாறு வரையறுக்கப்படும் $h\colon \Nat \to
%\Bin^\omega$ ஐக் கருதுக:
%\[
%h(n) = \underbrace{000\dots0}_{\text{$n$ $0$'s}}
%\]
%இது ஒரு சார்பு; ஆனால் $\Nat$ என்பது !!{enumerable}.
%\end{explain}

% SEGMENT OLP-0040-S06
% LOCAL SOURCE CORRECTION TA-SIZ-LOCAL-003; see translation/size-notes.tex.
\begin{prob}\label{sfr:siz:red-alt:prob:nat-nat}
  $f\colon \Nat \to \Nat$ என்ற எல்லாச் சார்புகளின் கணம் $X$ என்பது
  !!{nonenumerable} என்பதை ஒரு குறைத்தல் வாதத்தின் மூலம் காட்டுக
  (குறிப்பு: $X$ இலிருந்து $\Bin^\omega$ க்குச் செல்லும்
  ஒரு மேற்கோர்த்தல் சார்பைத் தருக).
\end{prob}

% SEGMENT OLP-0040-S07
\begin{prob}
இயல் எண் ஜோடிகளின் எல்லாக் \emph{கணங்களுடைய} கணம்,
அதாவது $\Pow{\Nat \times \Nat}$, !!{nonenumerable} என்பதை
ஒரு குறைத்தல் வாதத்தின் மூலம் காட்டுக.
\end{prob}

% SEGMENT OLP-0040-S08
\begin{prob}
இயல் எண்களின் முடிவுறாத தொடர்களுடைய கணம் $\Nat^\omega$ என்பது
!!{nonenumerable} என்பதை ஒரு குறைத்தல் வாதத்தின் மூலம் காட்டுக.
\end{prob}

% SEGMENT OLP-0040-S09
%\begin{prob}
%$P$ என்பது $\Nat$ இலிருந்து $\{0\}$ க்குச் செல்லும் சார்புகளின் கணம் என்க;
%$Q$ என்பது நேர்ம முழுக்களின் கணத்திலிருந்து $\{0\}$ க்குச் செல்லும்
%\emph{பகுதி}ச் சார்புகளின் கணம் என்க.
%$P$ என்பது !!{enumerable} என்றும் $Q$ அவ்வாறல்ல என்றும் காட்டுக.
%(குறிப்பு: $\Bin^\omega$ ஐப் பட்டியலாக்கும் சிக்கலை,
%$Q$ ஐப் பட்டியலாக்கும் சிக்கலாகக் குறைக்கவும்.)
%\end{prob}

% SEGMENT OLP-0040-S10
\begin{prob}
$S$ என்பது $\Nat$ இலிருந்து $\{0,1\}$ என்ற கணத்திற்குச் செல்லும்
எல்லா !!{surjection}s இன் கணம் என்க;
அதாவது, $S$ என்பது எல்லா !!{surjection}s~$f \colon \Nat
\to \Bin$ ஐயும் கொண்டுள்ளது.
$S$ என்பது !!{nonenumerable} எனக் காட்டுக.
\end{prob}

% SEGMENT OLP-0040-S11
\begin{prob}
எல்லா மெய்யெண்களின் கணம் $\Real$ என்பது !!{nonenumerable} எனக் காட்டுக.
\end{prob}

\end{document}
